\section{Scope, source baseline, and principal results}
\label{intro:sec}

The problem is to turn specific numerical conjectures and qualitative
asymptotic questions in the ProveIt project into exact mathematics.
The source baseline is commit
\texttt{4c173c06cc32c9cea554b39be1837ad2ae897fc1}, including the
manuscript directory and all six archives then present in
\src{docs/incoming}~\cite{baseline}. The supplied
\src{data/source_manifest.json} records that commit, the archive hashes,
and the principal source files used here. An \emph{open question} below
means open in that inspected material, unless a broader literature
statement is explicitly supported.

The article concentrates on four targets with a direct connection to
that material: a rational tetralogarithm formula, fractional Euler
remainders and their extremal constants, comparison at fixed total
order, and uniform saturation of Lerch zero counts. This focus permits
full proofs and reproducible finite certificates. The higher Gaussian
$S_{2m}$ reductions and the higher golden ladders remain distinct
problems.

\subsection{From a numerical relation to an exact functional identity}

The rational weight-four formula in
\src{03-algebraic.tex}, label \src{golden:eq:tetra}, was proposed in
2015 and repeated in 2017~\cite{Reshetnikov2015,Reshetnikov2017}.
The manuscript reduces it to a particular cross-base identity, but
does not derive that identity with its required branch information.
Theorem~\ref{rattetra:thm:main} completes the proof.

The certificate has forty-two rational functions of one variable.
Their exact tensor cancellation implies constant identities for three
modified real polylogarithms. A descent through weights four, three,
and two then restores the ordinary tetralogarithms, with all endpoint
and inversion contributions retained. Theorem~\ref{rattetra:thm:functional}
states the resulting functional family for every $0<t<1$.
The original conjecture is its $t=1/2$ consequence, while
Corollary~\ref{rattetra:cor:onethird} gives a further explicit
twenty-five-argument identity at $t=1/3$.

The finite certificate is part of the proof: its complete coefficient
table is printed, and its exact integer-coordinate cancellation is
replayable. Numerical polylogarithm evaluations serve only to check
the branch conventions independently.

\subsection{Two parameter limits that require different arguments}

Let $g_{a,b}$ be the imaginary part at $i$ of the depth-two harmonic
polylogarithm, let $E_N(a,b)$ be its $N$-term Euler transform, and write
\[
 R_N(a,b)=2^N(E_N(a,b)-g_{a,b}),\qquad
 C_N=\sup_{a,b>0}R_N(a,b).
\]
Precise definitions are given in \eqref{frac:eq:defs}. For fixed
positive $a,b$, the new leading law is
\[
 E_N(a,b)-g_{a,b}\sim
 \frac{\sqrt\pi}{2^{a+b}\Gamma(a+b)}
 2^{-N}N^{-1/2}(\log N)^{a+b-1}.
\]
Theorem~\ref{frac:thm:allorders} supplies every logarithmic order,
uniformly on compact positive parameter sets. Its two coefficient
sequences explain both termination at integral parameters and the
factorial divergence in Proposition~\ref{frac:prop:divergence}.

Taking a supremum over the whole positive quadrant is a different
limit. The incoming uniform-bounds report proves $C_N\to1$ and proposes
the more precise rate~\cite{incoming-bounds}. We prove
\[
 C_N-1\sim cN^{-p},\qquad
 p=\frac{\log2}{\log(3/2)},\qquad
 c=(1-q^{-1})q^{-p},\qquad q=\frac{\log3}{\log2}.
\]
There is no contradiction with the fixed-parameter law: the maximizing
inner order grows as $\log N/\log(3/2)$, and the outer order tends to
zero. These parameters leave every compact positive set.

Theorem~\ref{lateaxis:thm:reduction} proves the exact reduction
\[
 C_N=\max_{b>0}R_N(0,b)
 \quad\text{for }N\ge\lceil e^{24}\rceil=26{,}489{,}122{,}130.
\]
That explicit threshold is sufficient and conservative.
Theorem~\ref{lateaxis:thm:sharpasymptotic} gives a second, strictly
smaller scale and the displacement of the maximizing inner order.
The strict monotonicity of the axis maxima holds already at every
$N\ge2$; its transfer to $C_N$ uses the eventual reduction.

\subsection{A complete fixed-total comparison}

At fixed $w=a+b$, the Gaussian value at $N=1$ is strictly below the
axis value. It is tempting to extend this ordering to every
truncation. The exact answer depends on the outer order.
For $a\ge1$, Theorem~\ref{phasefrac:thm:above} proves the strict
comparison at every $N$, together with a positive Hausdorff-moment
representation of the difference.
For $0<a<1$, Theorem~\ref{phasefrac:thm:unique} proves exactly one
simple reversal in the natural continuous interpolation of $N$.
For $a=b=1/2$, exact rational enclosures show
\[
 R_N(1/2,1/2)<R_N(0,1)\quad(1\le N\le20),\qquad
 R_N(1/2,1/2)>R_N(0,1)\quad(N\ge21).
\]
The infinite classification follows from the sign-change theorem
and two finite certified signs. It is not inferred from a plot.

\subsection{A polynomial sufficient order for all Lerch zeros}

The reciprocal-gamma Appell polynomials are defined by
\[
 \frac{e^{xz}}{\Gamma(1+z)}
   =\sum_{n\ge0}Q_n(x)\frac{z^n}{n!}.
\]
Theorem~\ref{polylerch:thm:mesh} proves
$\mesh(Q_n)\ge4/[n(n-1)(n-2)]$ for $n\ge3$.
Combined with a gamma concentration estimate that uses the distances
to all roots, this yields the sufficient derivative order
\[
 \mathcal K_n=
 \left\lceil
 \left(16+6n(n-1)(n-2)H_{n-1}\right)^2
 \right\rceil
 \sim36n^6\log^2 n.
\]
For every $k\ge\mathcal K_n$ and every $0\le\rho\le1$, the Lerch
derivative family has exactly $n$ positive zeros, all simple. If
$x_{n,1}>\cdots>x_{n,n}$ are the roots of $Q_n$, their locations obey
\[
 \left|\log\frac{a_{n,j,k}^{\rho}}{ke^{-x_{n,j}}}\right|
 <\frac8{\sqrt k}.
\]
This permits a simultaneous range
$n\le c\,k^{1/6}/(\log k)^{1/3}$ for any fixed $c<1$.
The classical mesh-preservation theorems are credited to their
sources~\cite{polylerch:KSV}; the new contribution is their explicit
application and the improved sign transfer.

\subsection{Conventions and the role of computation}

All logarithms of positive real numbers are natural logarithms.
For real polylogarithm arguments outside the unit interval,
$\ell_n(x)=\Rea\Li_n(x)$ denotes the real part of the principal
continuation. Reciprocal gamma factors are interpreted as entire
functions, including their zeros at nonpositive integers.
In the Euler sections $a$ is the outer order, $b$ the inner order,
and $N$ the truncation index. In the Lerch section $n$ is the
logarithmic index and $k$ the derivative order. The letter $\rho$
there ranges over the full closed interval $[0,1]$.

The proof of each infinite assertion is given in the text. The exact
programs verify finite tensor identities, symbolic normalizations,
rational root brackets, and sign enclosures. The floating-point
program evaluates positive integrals to illustrate the rate of
convergence; it does not certify the global supremum or an asymptotic
error bound. No proof-assistant formalization is claimed. The archive
contains the complete source, the generated PDF, the replay programs,
frozen outputs, a source manifest, and proposed integration notes.
